\section{Failure Modes and Recovery}
\label{sec:failure-modes}

This section catalogues the failure modes TARDUS is designed to survive,
specifies the recovery procedure for each, and documents the residual
classes that operators must monitor for. The catalogue is the operational
counterpart of \S7's adversary model: \S7 says \emph{which adversaries
the protocol resists}, \S8 says \emph{what happens when something actually
goes wrong} and \emph{how to detect and recover}.

Failure modes are grouped by origin: validator-side (\S8.1),
user-side (\S8.2), protocol-level adversarial (\S8.3), and
network-level (\S8.4). Each entry lists \textit{Trigger},
\textit{User-facing consequence}, \textit{Recovery}, and
\textit{Detection signal}.

\subsection{Validator-Side Failure Modes}
\label{sec:failure-validator}

\paragraph{F1 --- Single validator crash.}
\begin{description}[leftmargin=2em,style=nextline,itemsep=0pt]
  \item[Trigger] One of the $n$ mint operators is offline (hardware
    fault, network drop, process restart).
  \item[Consequence] If the remaining online operators number $\geq t$,
    the mint continues to issue, refresh, and rotate normally; the
    offline operator simply does not contribute to active sessions.
    If online $< t$, all new issuance and refresh sessions stall;
    \emph{withdrawals} continue without disruption (on-chain only).
  \item[Recovery] The crashed operator restarts, loads its persistent
    share state from encrypted disk (\S3.1 HSM-mediated storage),
    and rejoins the next ceremony. No share regeneration required.
  \item[Detection] Operator heartbeat absent from the transparency log
    for $>$ 30 s. The validator daemon (Faz 2) publishes a liveness
    signal every 10 s; three consecutive misses raise an operator
    on-call alert.
\end{description}

\paragraph{F2 --- Validator network partition.}
\begin{description}[leftmargin=2em,style=nextline,itemsep=0pt]
  \item[Trigger] A network split isolates a subset $P \subsetneq [n]$
    of operators from the rest.
  \item[Consequence] If $|P| < t$ and $|[n] \setminus P| \geq t$,
    the majority side continues; minority side stalls.
    If both sides are $\geq t$, neither side should sign autonomously
    (a brain-split risk for the joint key). The validator daemon's
    quorum check enforces this: each side independently observes
    the other side's absence and pauses.
  \item[Consequence (rare)] If both sides are $\geq t$ and both ignore
    the quorum check (operator misconfiguration), both sides can issue
    valid signatures under the same joint key, creating fork-style
    double-issuance. This is operationally catastrophic but the
    protocol still rejects on-chain double-spend (T3) — the conflict
    surfaces at the first user who tries to spend two coins that
    happen to share a nullifier ancestry. Recovery requires a manual
    reconciliation against the transparency log.
  \item[Recovery] Partition heals → next ceremony auto-resumes. If
    forked-issuance occurred, manual epoch rotation via a key-revoke
    + re-register cycle (\S5.3) under committee consensus.
  \item[Detection] Cross-side keepalive failure → operator on-call;
    automated quorum-check pause is the first line of defence.
\end{description}

\paragraph{F3 --- HSM failure.}
\begin{description}[leftmargin=2em,style=nextline,itemsep=0pt]
  \item[Trigger] An operator's FIPS 140-2 Level 3 HSM (\S3.1) hardware
    fails or the key wrap protocol detects tampering.
  \item[Consequence] That operator can no longer participate; treat as
    F1 (single-validator-crash) extended.
  \item[Recovery] Replace the HSM, re-initialise with a fresh share
    derived from a proactive reshare (\S3.7). The reshare protocol's
    zero-polynomial property (T5) means the joint key is preserved
    and the failed operator's old share becomes useless — even if
    later recovered from the failed HSM. Operator's share is replaced
    in-epoch; no user-facing impact beyond F1-class downtime.
  \item[Detection] HSM self-test failure; share key derivation
    refuses to produce a result.
\end{description}

\paragraph{F4 --- Single share material leak.}
\begin{description}[leftmargin=2em,style=nextline,itemsep=0pt]
  \item[Trigger] One operator's share leaves the HSM boundary
    (insider attack, side channel, memory dump).
  \item[Consequence] By T1, no signature forgery yet (leaked share
    counts toward the $t-1$ tolerated minority).
  \item[Recovery] Immediate proactive reshare (\S3.7); the leaked share
    becomes useless one epoch later. The transparency log records the
    rotation event for auditability.
  \item[Detection] HSM audit log shows export attempt; or threat intel
    feed flags the leaked material. This is the hardest mode to
    detect — the operator-side organisational controls (background
    checks, dual-control share access, off-site backup) are the
    primary mitigation.
\end{description}

\paragraph{F5 --- $t-1$ share leak.}
\begin{description}[leftmargin=2em,style=nextline,itemsep=0pt]
  \item[Trigger] $t-1$ operators' shares leak simultaneously (concerted
    insider attack, supply-chain compromise).
  \item[Consequence] Adversary at boundary of T1; needs one more
    share to forge.
  \item[Recovery] Emergency reshare, on-chain key revoke
    (\texttt{Revoke} instruction, \S5.3), and re-register with a fresh
    DKG. All in-flight issuance pauses; existing coins remain valid
    (their signature was produced before the leak).
  \item[Detection] Multiple HSM audit log anomalies in a window;
    correlation with operator personnel events.
\end{description}

\paragraph{F6 --- $t$ or more share leak (catastrophic).}
\begin{description}[leftmargin=2em,style=nextline,itemsep=0pt]
  \item[Trigger] $\geq t$ operators' shares leak (or $t$ operators
    actively collude).
  \item[Consequence] Joint secret $x$ reconstructable via Lagrange
    interpolation (\S3.6); adversary can forge arbitrary coins
    under that joint key. T1 no longer holds for the affected epoch.
  \item[Recovery] On-chain revocation of the compromised keyset
    (\texttt{Revoke} instruction freezes new spends), forensic audit,
    user-side claim against the vault collateral (\S8.6 outlines the
    governance process for slashing the compromised operators' bonds
    if applicable). \emph{No protocol-level rollback}: existing
    transactions in the chain are not reversed.
  \item[Detection] Anomalous issuance volume; on-chain analytics flag
    a spike in newly registered nullifiers without matching
    pre-deposit transfers.
\end{description}

\subsection{User-Side Failure Modes}
\label{sec:failure-user}

\paragraph{F7 --- Lost coin material.}
\begin{description}[leftmargin=2em,style=nextline,itemsep=0pt]
  \item[Trigger] User deletes or loses their wallet's coin record
    (the $(x, Cp, \sigma)$ triple).
  \item[Consequence] The coin's collateral becomes irretrievable
    \emph{unless} the user holds the AEAD-encrypted backup
    (\S6 wallet layer, \texttt{backup.rs}). TARDUS coins are bearer
    instruments: there is no protocol-level recovery from loss.
  \item[Recovery] Restore from \texttt{seal\_backup}/\texttt{open\_backup}
    using the BIP-39 mnemonic-derived master seed (\S6.4 wallet
    backup). If the user has no backup, the coin's lamports remain
    locked in the vault, unspendable, until the protocol's collateral
    expiry rule (operationally: never; the vault retains them
    indefinitely as overcollateralisation).
  \item[Detection] User-side only; the protocol has no view into
    wallet contents.
\end{description}

\paragraph{F8 --- Wallet device theft.}
\begin{description}[leftmargin=2em,style=nextline,itemsep=0pt]
  \item[Trigger] User's device is stolen with unencrypted coin store.
  \item[Consequence] Thief can spend any unspent coin (bearer model).
  \item[Recovery] User submits \texttt{Refresh} transactions for every
    coin they remember (race against the thief). The first to reach
    the on-chain nullifier set wins.
  \item[Detection] User-side only; consider hardware-wallet integration
    (Faz 3 wallet roadmap) for stronger device-level mitigation.
\end{description}

\paragraph{F9 --- Failed refresh session.}
\begin{description}[leftmargin=2em,style=nextline,itemsep=0pt]
  \item[Trigger] Network drop, validator partition (F2), or user
    timeout during the 6-round refresh protocol (\S4).
  \item[Consequence] If the session aborted \emph{before} Round 3
    (mint's challenge), the user retains the old coin intact and may
    retry. If aborted \emph{after} Round 5 (mint's signed reveal),
    the user holds an unfinalised new coin but the old coin is
    nullified on chain — the user must complete Round 6 (off-chain
    unblind) using the persisted refresh-session state from
    \texttt{tardus-client::backup.rs} to recover the new coin.
  \item[Recovery] The wallet persists per-session state on every round
    transition; on restart it reads the state and resumes from the
    last completed round.
  \item[Detection] User-facing wallet timeout; the session state
    encodes its own progress.
\end{description}

\paragraph{F10 --- Refresh session race (network reordering).}
\begin{description}[leftmargin=2em,style=nextline,itemsep=0pt]
  \item[Trigger] Two refresh sessions on the same coin (e.g.\ retry
    of a session whose status is unclear) reach the on-chain program
    in different orders across validators.
  \item[Consequence] First-to-finalise wins; the second aborts with
    \texttt{ERR\_DOUBLE\_SPEND} (T3). User-facing: they hold one valid
    new coin instead of two.
  \item[Recovery] None needed; this is the intended atomic behaviour.
    The CLI's preflight nullifier check (\texttt{tardus devnet refresh})
    surfaces the duplicate before the second TX pays a fee.
  \item[Detection] CLI preflight; or post-finalisation \texttt{Refresh}
    transaction error.
\end{description}

\subsection{Protocol-Level Adversarial Failures}
\label{sec:failure-adversarial}

\paragraph{F11 --- DKG ceremony cheater.}
\begin{description}[leftmargin=2em,style=nextline,itemsep=0pt]
  \item[Trigger] During DKG, one or more operators publish broadcasts
    inconsistent with the shares they send privately (deviating from
    \S3.5).
  \item[Consequence] Honest operators detect the inconsistency via the
    dual-commitment check ($C_{i,0} = A_{i,0}$ in canonical form
    over both bases $G$ and $H$). The ceremony aborts before any
    operator commits to a final share.
  \item[Recovery] Restart DKG with the cheater excluded; publish the
    cheating proof to the transparency log.
  \item[Detection] Automatic during DKG round 2 validation.
\end{description}

\paragraph{F12 --- Reshare cheater (non-zero polynomial).}
\begin{description}[leftmargin=2em,style=nextline,itemsep=0pt]
  \item[Trigger] During proactive reshare (\S3.7), an operator's
    delta-polynomial $\tilde{f}_i$ has $\tilde{f}_i(0) \neq 0$
    (which would shift the joint key).
  \item[Consequence] Honest operators detect via the
    $\tilde{A}_{i,0} = \mathcal{O}$ check (added in v1.4.3 spec
    discovery). The reshare aborts; old shares remain valid for
    the current epoch.
  \item[Recovery] Restart reshare excluding the cheater; operators'
    next-epoch keyset rotation is delayed by one cycle.
  \item[Detection] Automatic during reshare validation.
\end{description}

\paragraph{F13 --- On-chain instruction spam (DoS).}
\begin{description}[leftmargin=2em,style=nextline,itemsep=0pt]
  \item[Trigger] Adversary submits high-volume invalid \texttt{Refresh}
    or \texttt{Withdraw} transactions to exhaust validator compute.
  \item[Consequence] Solana's transaction fee market increases TX cost
    during congestion; honest users pay more but the protocol remains
    available (per-instruction CU is bounded at $\sim$73k, well within
    the 200k default budget, so an attacker cannot cheap-spam).
  \item[Recovery] Solana's fee dynamics auto-rate-limit; no protocol
    action required.
  \item[Detection] Standard Solana RPC metrics.
\end{description}

\paragraph{F14 --- Nullifier-set capacity exhaustion.}
\begin{description}[leftmargin=2em,style=nextline,itemsep=0pt]
  \item[Trigger] The Borsh-encoded \texttt{NullifierSet} PDA reaches its
    8\,KiB allocation cap ($\approx 250$ entries).
  \item[Consequence] New \texttt{Refresh} and \texttt{Withdraw}
    transactions fail with \texttt{ERR\_ACCOUNT\_DATA\_TOO\_SMALL}.
  \item[Recovery] v1.4.13 migration to Light Protocol's compressed
    Merkle tree raises capacity to $2^{32}$ leaves with constant
    on-chain storage (root only). Until then, operators monitor
    nullifier count and reallocate the PDA (a one-time event for
    the lifetime of the deployment, since the BPF Loader Upgradeable
    permits account-data reallocation by the program authority).
  \item[Detection] \texttt{tardus devnet info} reports nullifier count;
    threshold-based alerting at 80\% of capacity.
\end{description}

\subsection{Network-Level Failures}
\label{sec:failure-network}

\paragraph{F15 --- Solana network outage.}
\begin{description}[leftmargin=2em,style=nextline,itemsep=0pt]
  \item[Trigger] Solana mainnet halts (block production stops).
  \item[Consequence] No new \texttt{Refresh} or \texttt{Withdraw}
    transactions land; existing coins remain valid; the mint
    committee continues to issue off-chain (issuance does not
    require on-chain TX). Withdrawals stall.
  \item[Recovery] Resume when Solana production resumes; pending
    transactions are picked up automatically by the validator
    daemon's retry queue.
  \item[Detection] Solana RPC reports stale slot height; transparency
    log records the pause.
\end{description}

\paragraph{F16 --- Solana chain reorganisation past finality.}
\begin{description}[leftmargin=2em,style=nextline,itemsep=0pt]
  \item[Trigger] An unprecedented Solana consensus event reorganises
    blocks at depth $\geq 32$ (the safety boundary in A6).
  \item[Consequence] Catastrophic for any Solana program; TARDUS
    inherits the network's safety property. T3 and T6 fail if A6
    fails.
  \item[Recovery] Coordinated with Solana ecosystem; nothing TARDUS
    can do unilaterally.
  \item[Detection] Solana RPC reports finalised-block divergence;
    network-wide alert.
\end{description}

\paragraph{F17 --- Token-2022 program bug (post-v1.4.12).}
\begin{description}[leftmargin=2em,style=nextline,itemsep=0pt]
  \item[Trigger] A bug in Solana's Token-2022 program affects the
    Confidential Mint extension (v1.4.12+).
  \item[Consequence] Vault SPL operations may revert or, worse, leak
    amounts that should be encrypted. v1.4.12 will document the
    explicit dependence on Token-2022 invariants.
  \item[Recovery] Coordinate with Solana Labs; pause vault operations
    via \texttt{Revoke} until the upstream fix is deployed.
  \item[Detection] Standard SPL ecosystem alerting.
\end{description}

\paragraph{F18 --- Malicious RPC node.}
\begin{description}[leftmargin=2em,style=nextline,itemsep=0pt]
  \item[Trigger] User's RPC provider returns stale or fabricated
    account data.
  \item[Consequence] CLI preflight check (\texttt{query-nullifier})
    could miss an already-inserted nullifier, leading to a wasted
    TX fee when the on-chain check catches it. The on-chain check
    always wins (T3), so no double-spend risk.
  \item[Recovery] Multi-RPC verification (Faz 3 wallet roadmap); or
    user-side fallback to the validator committee's own RPC.
  \item[Detection] Cross-RPC discrepancy alarms; user-facing warning.
\end{description}

\subsection{Relay-Layer Failure Modes}
\label{sec:failure-relay}

The relay layer (\Cref{sec:relay}) introduces its own operational
F-modes catalogued as R1--R7 in \Cref{sec:relay-failure-modes}.
Cross-referenced summary:

\paragraph{R1 --- Relay process crash.} F1-class extension: with
v5.4 \texttt{SQLite} persistence (\Cref{sec:relay-storage}) the
in-flight inbox survives restart; the validator-side equivalent of
F1's ``operator restart'' applies directly.

\paragraph{R2 --- Relay disk full.} F1-class extension: \texttt{SQLite}
writes fail $\to$ POST returns HTTP 5xx $\to$ sender's wallet
retries with backoff. No coin-side state damage.

\paragraph{R3 --- Relay operator payload inspection.} Mitigated by
the v5.6 sealed-box payload format (\Cref{sec:relay-payload}):
the relay cannot recover coin material under T7. A relay-side audit
still sees recipient pubkey, timestamps, ciphertext sizes (traffic
metadata) — Tor/I2P circuiting is the deployment-time defence.

\paragraph{R4 --- Relay log seizure.} Same mitigation as R3: logs
hold only opaque ciphertext + metadata.

\paragraph{R5 --- Inbox exhaustion attack.} F13-class extension: a
malicious sender deposits up to \texttt{max\_per\_recipient} entries
on the victim's inbox. The v5.6 mitigation is the per-recipient cap
(503 on overflow) + TTL-based eviction. Long-term mitigation is a
fee market (deferred to v5.7).

\paragraph{R6 --- Recipient secret leaked.} F8-class extension to
the wallet section: leak of the recipient's \texttt{recv\_sk}
(\Cref{sec:wallet-recv-identity}) lets the attacker decrypt all
ciphertexts addressed to that pubkey, past and future. The wallet
mnemonic is therefore the security root for both coin custody and
P2P payload privacy.

\paragraph{R7 --- Sender IP fingerprinting.} Not modelled in v5; a
relay-side packet sniffer can link a recipient pubkey to a sender IP.
Production deployments are expected to circuit sender wallets
through Tor or a privacy-respecting VPN before reaching the relay.

\subsection{HSM-Backend Failure Modes}
\label{sec:failure-hsm}

The v2.11 PKCS\#11 backend
(\Cref{sec:mint-validator-daemon}) introduces failure modes that the
v2.1 file backend did not have. These extend F3 (HSM failure):

\paragraph{F3a --- PKCS\#11 module path missing.}
\begin{description}[leftmargin=2em,style=nextline,itemsep=0pt]
  \item[Trigger] \texttt{--hsm-module} points to a missing or
    incompatible \texttt{.so}.
  \item[Consequence] Validator boot fails with explicit
    \texttt{Error::Config}; no shares ever loaded.
  \item[Recovery] Fix path; rotate operator if the wrong module was
    deployed.
\end{description}

\paragraph{F3b --- PKCS\#11 user PIN rotated upstream.}
\begin{description}[leftmargin=2em,style=nextline,itemsep=0pt]
  \item[Trigger] HSM operator rotates the user PIN through the
    HSM-vendor management console without updating the validator's
    config.
  \item[Consequence] Validator boot fails at the login step
    (\texttt{C\_Login} returns \texttt{CKR\_PIN\_INCORRECT}); no
    silent fallback.
  \item[Recovery] Sync new PIN; validator restart.
\end{description}

\paragraph{F3c --- Wrap key deleted on the HSM.}
\begin{description}[leftmargin=2em,style=nextline,itemsep=0pt]
  \item[Trigger] Operator misclick or HSM tenancy event removes the
    AES wrap key (\texttt{CKA\_LABEL = tardus-wrap-v1}).
  \item[Consequence] All on-disk \texttt{share\_*.bin} files become
    permanently undecryptable; the affected validator must
    re-participate in a reshare to obtain a fresh share. This is
    F4-class for the affected validator but does not compromise the
    threshold.
  \item[Recovery] Reshare under existing \texttt{joint\_pk}; the
    re-keyed validator picks up its new share from the reshare
    transcript.
\end{description}

\paragraph{F3d --- HSM clock or session timeout.}
\begin{description}[leftmargin=2em,style=nextline,itemsep=0pt]
  \item[Trigger] HSM enforces session lifetime; the
    \texttt{Pkcs11ShareStore}'s long-held session is invalidated.
  \item[Consequence] Next \texttt{C\_Encrypt}/\texttt{C\_Decrypt}
    fails with \texttt{CKR\_SESSION\_HANDLE\_INVALID}.
  \item[Recovery] v2.12 will add automatic session re-open on
    timeout; v2.11 requires validator restart.
\end{description}

\subsection{Detection and Monitoring Discipline}
\label{sec:failure-monitoring}

The full failure-mode catalogue above is operationalised through three
observation channels:

\begin{enumerate}[leftmargin=2em,itemsep=0pt]
  \item \textbf{Transparency log.} Append-only, hash-chained log
    published by the mint committee, recording every DKG round,
    every reshare event, every keyset registration/revocation, and
    every operator heartbeat. Anyone can replay the log to verify
    committee behaviour. F1, F2, F4, F5, F11, F12 surface here.
  \item \textbf{On-chain analytics.} Standard Solana indexers query
    the program's account states. \texttt{tardus devnet info}
    (\S6.5) is the minimal CLI front for this; richer dashboards
    (Faz 3) will track nullifier-rate, vault-balance, and
    keyset-registry growth. F6, F13, F14, F18 surface here.
  \item \textbf{Operator alerting.} Per-validator daemon (Faz 2)
    Prometheus-style metrics on heartbeat, share count, ceremony
    participation, HSM health. F1, F2, F3, F4 surface here first.
\end{enumerate}

\subsection{Governance and Open Problems}
\label{sec:failure-open}

\paragraph{Slashing.} The current v1 protocol does not implement
operator-bond slashing; F4-F6 recovery depends on out-of-band
governance (e.g.\ legal action against an operator who leaked
material). A v2 extension under consideration: operators post a
collateral bond into a separate Solana PDA; the
\texttt{Revoke} instruction can be augmented with a slashing
parameter that diverts the bond to the vault as compensation. This is
deferred to Faz 5.

\paragraph{Cross-chain bridge attacks.} Coins denominated in
bridged assets (e.g.\ Wormhole-wrapped tokens) inherit the bridge's
security posture. TARDUS itself cannot defend against a bridge
compromise. This will be addressed in the Faz 5 mainnet asset-class
policy: only assets with native Solana issuance (SOL, Token-2022 mints
with the issuer's signature attestation) qualify for the canonical
TARDUS keyset; bridged assets get separate keysets with explicit
risk disclosure.

\paragraph{Post-quantum migration.} A1 (ECDLP) is broken by a
sufficiently large quantum computer (Shor's algorithm, $\sim 4000$
logical qubits). TARDUS will need a post-quantum signature primitive
in that era; current candidates (Dilithium, Falcon) have $\sim 2$\,KB
signatures, incompatible with the per-coin storage budget. The
mitigation will be a parallel post-quantum keyset (separate
\texttt{KeysetEntry} with a post-quantum algorithm marker), with the
classical keyset deprecated over a multi-year migration window. See
\S5 for the keyset registry's extensibility.

\paragraph{Async secure messaging.} Faz 2's validator daemon needs
authenticated, replay-protected, end-to-end-encrypted messaging
between operators. The current design assumes TLS+mTLS over a
private network; a fully decentralised alternative
(libp2p, Iroh) is under evaluation.

\subsection{Summary Table}
\label{sec:failure-summary}

\begin{center}
\small
\begin{tabular}{@{}llll@{}}
\toprule
ID & Mode & Severity & Recovery class \\
\midrule
F1  & Single validator crash       & Low      & Operator restart (v2.X)         \\
F2  & Validator network partition  & Medium   & Auto-quorum-pause; manual recon \\
F3  & HSM failure                  & Low      & Reshare + replace               \\
F4  & Single share leak            & Medium   & Proactive reshare               \\
F5  & $t-1$ share leak             & High     & Emergency reshare + revoke      \\
F6  & $\geq t$ share leak          & Catastrophic & On-chain revoke + audit     \\
F7  & Lost coin material           & High (user) & AEAD backup restore          \\
F8  & Wallet theft                 & High (user) & Race-refresh; HW wallet      \\
F9  & Refresh session abort        & Low      & Wallet auto-resume              \\
F10 & Refresh race                 & None     & First-finalised wins (T3)       \\
F11 & DKG ceremony cheater         & Low      & Auto-detect + restart           \\
F12 & Reshare cheater              & Low      & Auto-detect + restart           \\
F13 & Instruction spam (DoS)       & Low      & Solana fee market               \\
F14 & Nullifier capacity exhausted & Medium   & Reallocation; v1.4.13 Light tree \\
F15 & Solana network outage        & Low      & Wait + auto-retry               \\
F16 & Solana finality break        & Catastrophic & Ecosystem-wide              \\
F17 & Token-2022 bug (v1.4.12+)    & Variable & Pause + upstream fix            \\
F18 & Malicious RPC                & Low      & Multi-RPC verification          \\
\bottomrule
\end{tabular}
\end{center}

Of the 18 catalogued modes, 14 have automated detection and protocol-level
recovery (F1-F3, F9-F18); 4 require human-in-the-loop response (F4-F8).
This is the explicit operational surface the validator daemon (Faz 2)
must instrument.
